% Section 2: Reinforcement learning for forecasting
\subsection{Reinforcement learning for forecasting}
\label{sec:rl_for_forecasting}

\subsubsection{Reinforcement learning as an ensemble combiner}
\label{sec:rl_ensemble}

The simplest way to put RL near forecasting is to treat ensemble model selection as
a sequential decision problem. \citet{Fu2022} formulate the combination of $K$ base
forecasters as a Markov decision process. The state is the recent history of the
series plus the recent performance of each base model, the action is the vector of
combination weights, and the reward is the negative one-step forecast error. They
train a deterministic policy with DDPG and demonstrate gains over equal-weight and
performance-weighted ensembles on standard benchmarks. The methodological depth is
modest. The same construction can be read as an online learning problem with a
non-convex action space, and the regret guarantees one would obtain from a Hedge or
EXP3 algorithm are no worse than those obtained from DDPG. We mention this strand
mostly to note that it exists.

\subsubsection{Reinforcement learning with decision-relevant rewards}
\label{sec:rl_brier}

A second use of RL is to train a probabilistic forecaster directly on a proper
scoring rule that doubles as the downstream decision loss. \citet{Turtel2025} fine-tune
a 14B-parameter reasoning model on a corpus of $10^4$ Polymarket prediction questions
augmented with $10^5$ synthetic questions, using on-policy policy gradient with the
Brier score serving as the reward signal. The training procedure adapts GRPO by
removing the per-question variance normalisation that would otherwise wash out the
calibration gradient, and uses a baseline-subtracted advantage of the ReMax kind. On
a held-out test set of $1{,}265$ questions, the seven-run ensemble reaches a Brier
score of $0.190$ with $95\%$ confidence interval $[0.178, 0.203]$, an expected
calibration error of $0.062$, and a simulated trading profit of $\$52$ against $\$39$
for OpenAI o1 under the same trading rule. The headline lesson, in their words, is
that RL trained with a proper scoring rule as the reward improves calibration and
turns calibration into profit in a real market.

\subsubsection{Approximate dynamic programming for forecasting}
\label{sec:adp_forecasting}

Sequential estimation has a longer pedigree in the dynamic programming tradition.
\citet{BertsekasRollout2022} provides the unifying frame. A measurement policy at
time $t$ is a function from the current posterior over an unknown $\theta$ to the
next observation point, and rollout improves a base policy by performing a one-step
lookahead under Monte Carlo simulation of the residual horizon. The certainty
equivalence approximation replaces the stochastic continuation cost with its
expectation conditional on the realised first step, which lowers the computational
burden without forfeiting the policy improvement guarantee whenever the base policy
already satisfies a mild contractivity condition. The framework covers Bayesian
optimisation, optimal stopping for sequential testing, and adaptive Bayesian
experiment design, and it returns the same vocabulary as standard RL.

\citet{Fernandes2024} applies the neuro-dynamic programming variant of this idea to
the gross domestic product nowcasting problem. The reward is not the cross-sectional
forecast error. It is the adjustment cost between successive predictions plus a
terminal accuracy penalty, which Fernandes argues better reflects the way a central
bank or a finance ministry uses a nowcast. The value function is approximated by a
single-hidden-layer feedforward network with $32$ units, trained by temporal
difference learning on quarterly data from $26$ European Union countries with
Portugal held out for evaluation. Over the period 2015Q1 to 2023Q4, the neuro-dynamic
programming nowcaster lowers the out-of-sample mean absolute error from $6.6$ to
$4.5$ percentage points of GDP range, a $32\%$ reduction over OLS, with the root
mean squared error falling from $8.3$ to $5.9$ percentage points of range, a $28\%$
reduction. The most striking finding is that training on foreign data transfers well
to Portugal, which alleviates the small-T problem that plagues macroeconomic
forecasting.

\citet{BenAmeur2025} extend the dynamic programming side of the literature into the
online and non-stationary regime. They give the first sublinear regret bounds for
online decision-focused learning. The DF-FTPL algorithm achieves static regret of
order $T^{-1/4}$ against the best fixed parameter in hindsight, and the DF-OGD
algorithm achieves dynamic regret of order $(1 + P_T)^{1/4} T^{-1/4}$, where $P_T$
is the path variation of the comparator. Both bounds are independent of the
parameter dimension up to a $\log \log d$ factor, which matters for high-dimensional
forecasting models.

\subsubsection{Decision-focused learning}
\label{sec:dfl}

The deepest sub-stream is decision-focused learning, which trains the forecast
model under a loss that reflects the downstream decision regret rather than the
prediction error. The framework was formalised by \citet{ElmachtoubGrigas2022} under
the name Smart Predict-then-Optimize, abbreviated SPO. Given features $x$ and a
linear cost vector $c$ that determines the objective $c^\top z$ of a downstream
optimisation problem $\min_{z \in \mathcal Z} c^\top z$, the SPO loss measures the
suboptimality of the decision induced by the predicted cost $\hat c(x)$, namely
$\mathrm{SPO}(\hat c, c) = c^\top z^\star(\hat c) - c^\top z^\star(c)$, where
$z^\star(c)$ is the optimal solution under the true cost. The SPO loss is
discontinuous in $\hat c$, so \citet{ElmachtoubGrigas2022} construct the convex
surrogate $\mathrm{SPO}^+(\hat c, c)$. Their main theorem establishes Fisher
consistency.

\begin{theorem}[\citealp{ElmachtoubGrigas2022}, Theorem 2]
\label{thm:spo_fisher}
Suppose the noise around the conditional mean cost $\mathbb E[c \mid x]$ is centrally
symmetric and the feasible region $\mathcal Z$ is a convex polyhedron. Then the
population minimiser of the SPO$^+$ loss is the conditional mean cost function. In
particular, the regression function $x \mapsto \mathbb E[c \mid x]$ minimises both
the Bayes risk of the SPO loss and the Bayes risk of the SPO$^+$ surrogate.
\end{theorem}

The Fisher consistency result is non-trivial, because the SPO$^+$ loss is not a
proper scoring rule in the sense of \citet{GneitingRaftery2007}. Symmetric noise is
needed to align the surrogate's minimiser with the SPO minimiser, and asymmetric
noise generically breaks the equivalence.

\citet{LiuGrigas2021} convert Fisher consistency into a finite-sample statement. They
prove that on convex polyhedra the SPO$^+$ calibration function $\delta(\varepsilon)$
satisfies $\delta(\varepsilon) \ge c \varepsilon^2$, which combined with a standard
Rademacher complexity bound on the surrogate excess risk yields an excess SPO risk
of order $n^{-1/4}$ for empirical risk minimisation over a Rademacher-bounded
hypothesis class. For strongly convex level sets the rate improves to $n^{-1/2}$,
matching the calibration of ordinary squared error.

\citet{DontiAmosKolter2017} operationalised the SPO programme for neural networks by
differentiating through the Karush-Kuhn-Tucker conditions of the downstream
optimisation. They apply task-based learning to three economic problems, namely
electricity grid scheduling on PJM data, energy storage arbitrage, and an inventory
stock balancing problem, and report regret reductions of $10\%$ to $40\%$ over models
trained with standard maximum likelihood and plugged into the same optimisation
routine. The implicit differentiation step is the methodological bridge from
batch DFL to deep learning at scale, and it is the building block on which most
subsequent DFL papers rest.

The DFL story is not uniformly positive. \citet{HuKallus2020} prove a counterintuitive
result for contextual linear optimisation under low near-dual-degeneracy. Under their
margin condition, the regret of a naive plug-in estimator that fits the conditional
mean by ordinary least squares decays at rate $n^{-1}$, while an integrated decision-aware
estimator achieves $n^{-2/3}$. The naive method exploits the curvature of the loss
near the decision boundary in a way that a decision-aware method does not. Their
result is a useful corrective. Decision-focus is a worthwhile asymptotic objective
only when the hypothesis class cannot recover the true conditional cost, which is
the realistic regime for economic forecasting but is not the same as the universal
claim that DFL dominates.

The modern frontier addresses misspecification head on. \citet{HuangGupta2024}
introduce the family of Perturbation Gradient losses. The PG loss is a Lipschitz
continuous difference-of-concave surrogate whose stochastic gradient approximates
the directional derivative of the plug-in decision loss via zeroth-order
estimation. The key theorem states that under a sub-Gaussian noise model and a
Rademacher-complexity bound $R_n$ on the hypothesis class, the population
minimiser of the PG loss attains a regret of order
$\widetilde O\bigl((dp/n)^{1/4}\bigr)$ relative to the best-in-class policy. The
guarantee is best-in-class rather than full-information optimal, but it holds under
misspecification, which is precisely the case where SPO$^+$ loses its Fisher
consistency property. Empirically, PG outperforms SPO$^+$ on synthetic and
real-data instances where the cost function does not lie in the hypothesis class.

\citet{Mandi2024} survey the broader landscape. They organise gradient-based
methods, comprising surrogate-loss methods like SPO$^+$ and implicit-differentiation
methods like that of \citet{DontiAmosKolter2017}, alongside gradient-free methods such
as smoothing and Bayesian approaches. They benchmark $11$ DFL techniques on $7$
canonical problems, including shortest path, knapsack, portfolio, travelling
salesman, and bipartite matching, and report that DFL is most useful precisely when
the underlying prediction problem is misspecified. On well-specified problems,
the gap between predict-then-optimise and decision-focused learning shrinks to
within statistical error.
